\section{\secsumproofs}
\label{sec_sum_proofs}

We have observed patterns and conjectured explicit formulas for sums and products in \Cref{sec_sum_conj}, including sums involving Fibonacci numbers, and in \Cref{sec_arith_triangle}, including sums involving binomial coefficients.  In a few special cases, we saw how to use the Binomial Theorem to prove summation formulas, but how do we prove the Binomial Theorem itself as well as other formulas?  In this section, we look at two proof techniques, including new special cases of mathematical induction and an extension of combinatorial proofs to summations.  We close by looking at special sums, such as the handshake formula from \Cref{thm_hs} and finite geometric sums.

%%%%%%%%%%%%%%%%%%%%%%%%%

\newcommand{\subpffsumproofsind}{Proof Format: Special Case of Mathematical Induction for Proving Sum and Product Formulas}
\subsection{\subpffsumproofsind}
\label{sub_pff_sum_proofs_induction}

We can use mathematical induction to prove explicit formulas for summations and products. We begin with the proof format for a summation.

\begin{pff}[Mathematical induction: special case proving an explicit formula for a sum]
\label{pff_math_ind_sum}
    Given the explicitly-defined sequence $a_k$ for $k \ge s$, prove that $\displaystyle \sum_{k=s}^na_k=$ (explicit rule) for $n \ge s$.
    \begin{proof}
        Write $f(n) = $(explicit rule).  Induct on $n$.

        First, when $n=s$ we know that $\displaystyle \sum_{k=s}^sa_k=a_s= \cdots $ and also $f(s)=$ (evaluate and simplify to get the same answer).

        Next, assume $\displaystyle \sum_{k=s}^{n-1}a_k=f(n-1)= $ (evaluate and simplify), for some integer $n > s$.  Then 
            $$\sum_{k=s}^na_k 
            = \left(\sum_{k=s}^{n-1}a_k\right) + a_n
            = f(n-1) + a_n  
            = \cdots \text{(use algebra)} \cdots 
            = f(n).$$
    \end{proof}
\end{pff}

The key first step in the last line of the proof uses the recursive definitions of partial sums, \Cref{rem_rec_def_partial_sum}.

Similarly, here is the proof format for a product.

\begin{pff}[Mathematical induction: special case proving an explicit formula for a product]
\label{pff_math_ind_prod}
    Given the explicitly-defined sequence $a_k$ for $k \ge s$, prove that $\displaystyle \prod_{k=s}^na_k=$ (explicit rule) for $n \ge s$.
    \begin{proof}
        Write $f(n) = $(explicit rule).  Induct on $n$.

        First, when $n=s$ we know that $\displaystyle \prod_{k=s}^sa_k=a_s= \cdots $ and also $f(s)=$ (evaluate and simplify to get the same answer).

        Next, assume $\displaystyle \prod_{k=s}^{n-1}a_k=f(n-1)= $ (evaluate and simplify), for some integer $n > s$.  Then 
            $$ \prod_{k=s}^na_k 
            = \left(\prod_{k=s}^{n-1}a_k\right) \cdot a_n
            = f(n-1) \cdot a_n 
            = \cdots \text{(use algebra)} \cdots 
            = f(n).$$
    \end{proof}
\end{pff}

The key first step in the last line of the proof uses the recursive definitions of partial products, \Cref{rem_rec_def_partial_prod}.

As usual, when we introduce a new proof format, we give you examples to copy and edit.

\begin{act}[Proving sum and product formulas using mathematical induction]
\label{act_math_ind_sum_prod_formulas}  
~
    \begin{actpart}
        \item \label{part_proof_ind_sum_powers2} Copy the following proof.

            Use mathematical induction to prove that $\displaystyle \sum_{k=0}^n2^k=2^{n+1}-1$ for $n \ge 0$.  

            \begin{quote} 
                \emph{Proof} Write $f(n) = 2^{n+1}-1$.  Induct on $n$.
            
                First, when $n=0$ we know that $\displaystyle \sum_{k=0}^02^k=2^0=1$ and also $ f(0)=2^{0+1}-1=2-1=1.$
            
                Next, assume $\displaystyle \sum_{k=0}^{n-1}2^k=f(n-1)= 2^{(n-1)+1}-1=2^n-1$ for some integer $n > s$. Then 
                    \begin{align*}
                        \sum_{k=0}^n2^k & = \left(\sum_{k=0}^{n-1}2^k\right) + 2^n = f(n-1) + 2^n = \left(2^n-1\right)+2^n \\ 
                        & = 2^n+2^n-1=2\cdot2^n-1=2^{n+1}-1 = f(n).
                    \end{align*}
            \end{quote}
            
        \item Edit the proof in (\ref{part_proof_ind_sum_powers2}) to prove that $\displaystyle \sum_{k=0}^n2 \cdot 3^k = 3^{n+1}-1$ for $n\ge 0$.  

        \item \label{part_proof_ind_prod_recip} Copy the following proof.

            Use mathematical induction to prove that $\displaystyle \prod_{k=1}^nk=n!$ for $n \ge 1$.  
       
            \begin{quote} 
                \emph{Proof} Write $f(n) = n!$.  Induct on $n$.
            
                First, when $n=1$ we know that $\displaystyle \prod_{k=1}^1k=1=1$ and also $ f(1)=1!=1.$
            
                Next, assume $\displaystyle \prod_{k=1}^{n-1}k=f(n-1)= (n-1)!$ for some integer $n > s$. Then 
                    $$\prod_{k=1}^nk 
                    = \left(\prod_{k=1}^{n-1}k\right) \cdot n
                    = f(n-1) \cdot n  
                    = \left((n-1)!\right) \cdot n
                    = n(n-1)!
                    =n!
                    = f(n).$$
            \end{quote}
            
        \item Edit the proof in (\ref{part_proof_ind_prod_recip}) to prove that $\displaystyle \prod_{j=1}^n \frac{1}{j} = \frac{1}{n!}$ for $n\ge1$.

        \item \label{part_proof_ind_sum_Fib_evens} Copy the following proof.

        Use mathematical induction to prove that $\displaystyle \sum_{k=1}^nF_{2k}=F_{2n+1}-1$  for $n \ge 1$ where, as usual, the Fibonacci numbers are defined by $F_0=0$, $F_1=1$, and $F_n=F_{n-1} + F_{n-2}$ for $n \ge 2$.

            \begin{quote} 
                \emph{Proof} Write $f(n) = F_{2n+1}-1$.  Induct on $n$.

                First, when $n=1$ we have $\displaystyle\sum_{k=1}^1F_{2k}=F_{2(1)}=F_2=1$ and also $$f(1) = F_{2(1)+1}-1=F_3-1=2-1=1.$$
            
                Next, assume    
                    $$\sum_{k=1}^{n-1}F_{2k} = f(n-1) = F_{2(n-1)+1}-1 = F_{2n-1}-1$$
                for some integer $n > 1$.
                Then
                    \begin{align*}
                        \sum_{k=1}^nF_{2k} & = \left( \sum_{k=1}^{n-1}F_{2k} \right) + F_{2n} 
                        = f(n-1) + F_{2n}  
                        = \left(F_{2n-1}-1\right)+F_{2n} \\
                        & = \left(F_{2n} + F_{2n-1}\right) -1 = F_{2n+1}-1 = f(n)
                    \end{align*}
                The next-to-last equality uses the recursive definition the Fibonacci numbers to conclude that the sum of the consecutive terms $F_{2n}$ and $F_{2n-1}$ equals the next term, $F_{2n+1}.$  
            \end{quote}
                
            Note: We saw this summation in \Cref{act_fib_sums_conj} (\ref{part_fib_sums_even_index}) and \Cref{exam_fib_sum_conj_even_index}.

            \item Edit the proof from (\ref{part_proof_ind_sum_Fib_evens}) to prove that $\displaystyle \sum_{k=1}^{n}F_{2k-1} = F_{2n}$ for $n \ge 1$. 
            
            Note: We saw this summation in  \Cref{sec_sum_conj} \Cref{exer_fib_sum_odd_terms_conj}. 
    \end{actpart}
\end{act}

%%%%%%%%%%%%%%%%%%%%%%%%%%%%%%%%%%%%%%%%%%
\newcommand{\subpffcombsums}{Proof Format: Combinatorial Proof of Sums}
\subsection{\subpffcombsums}
\label{sub_pff_comb_sums}

We can use mathematical induction to prove the explicit formula of a sum of the form $\displaystyle \sum_{k=s}^na_k=f(n)$, but the recursive step of these proofs is based on the fact that the terms $a_k$ do not depend on the upper limit $n$. How can we prove an explicit formula for a summation where the terms depend on $n$? For example, in \Cref{sec_arith_triangle} we saw sums with binomial coefficients where the terms involved $n$. An option is to extend the proof format for a combinatorial proof, \Cref{pff_comb_proof}, to apply to sums.  This extended proof format is especially useful for proofs involving binomial coefficients, since, by definition, they count objects.

\begin{pff}[Combinatorial proof of sums]
\label{pff_comb_proof_sum}  
    Given the explicitly-defined sequence $a_k$ for $k \ge s$ where the terms $a_k$ may depend on $n$, prove that $\displaystyle \sum_{k=s}^na_k=f(n)$ for $n \ge s$.
        \begin{proof}
            We use a combinatorial proof. Consider the following situation \ldots.

            How many ways are there to \ldots?

            First, condition on $k=$ (say what $k$ counts).  In this case, there are (show how to get $a_k$) ways to do \ldots.  Since cases add (\Cref{thm_cases_add}), there are a total of $\displaystyle \sum_{k=s}^na_k$ ways.

            On the other hand, (explain how to count using a different method to get $f(n)$.)

            Since we counted the same quantity in two different ways, it follows that $\displaystyle \sum_{k=s}^na_k=f(n)$.
        \end{proof}
\end{pff}

The phrase ``condition on $k=$ (say what $k$ counts)'' indicates that we are considering a list of cases Case $s$, Case $s+1$, \ldots, Case $n$ where in Case $k$, there are $k$ of some object. Let's see this new proof format in action.

\begin{act}[Combinatorial proof of sum]
\label{act_comb_pf_sum}
~  
    \begin{actpart}
        \item \label{part_comb_pf_sum_row_sum_arith_tri} Copy the following proof.

            Use a combinatorial proof to prove that $$\displaystyle \sum_{k=0}^n\binom{n}{k} =2^n$$ for any natural number $n$. We made this conjecture about the row sums of the arithmetic triangle in \Cref{act_pascals_patterns} (\ref{part_find_row_sum_arith_tri}).
        
            \begin{quote} 
                \emph{Proof.}  We use a combinatorial proof.  Consider the following situation: we have the set $T = \{1,2,\ldots, n\}$.  How many ways are there to choose a subset of $T$?  

                First, condition on $k=$ the number of elements in the subset.  There are $\binom{n}{k}$ subsets with $k$ elements for $k = 0, 1, 2, \ldots, n$. Since cases add (\Cref{thm_cases_add}), there are a total of $\displaystyle \sum_{k=0}^n\binom{n}{k}$ subsets.  
        
                On the other hand, for each element of $T$ we can decide whether it will be in the subset or not (2 choices). There are two choices for whether 1 is in the subset, two choices for whether 2 is in the subset, and so on up to two choices for whether $n$ is in the subset.  Since steps multiply (\Cref{thm_steps_multiply}), there are $\underbrace{2\cdot 2 \cdots \cdots 2}_{n \text{ times}} = 2^n$ subsets.

                Since we counted the same quantity in two different ways, it follows that $$ \sum_{k=0}^n\binom{n}{k} =2^n.$$ 
            \end{quote}
      
        \item Edit the proof in (\ref{part_comb_pf_sum_row_sum_arith_tri}) to prove that $$\displaystyle \sum_{k=1}^nk\binom{n}{k} =n2^{n-1}$$ for any integer $n \ge 1$.  
        
        Hints: Out of a group of $n$ people, choose a subset to be a committee (with at least one person) where one of the people on the committee is the chair. Condition on $k =$ the number of people on the committee. In one way of counting, choose the committee members first and then the chair.  In the other way of counting, choose the chair first and then the rest of the committee members.      
    \end{actpart}
\end{act}

%%%%%%%%%%%%%%%%%%%%%%%%

\newcommand{\subspecialsum}{Detour: Special Sums}  
\subsection{\subspecialsum}
\label{sub_special_sums}

We have seen a variety of specific summations that have an explicit formula.  In this section, we look at several common classes of summations that have an explicit formula.  We begin with an activity.

\begin{act}[Special sums]
\label{act_special sums}
~
    \begin{apart}
        \item \label{part_sum_powers2_to4} Consider a real number $r \neq 1$ and the sum $$S =  \sum_{k=0}^4r^k = 1 + r + r^2 + r^3 + r^4.$$   Note that $S$ is a function of $r$. Evaluate $S$ when $r = 10$.
        
        \item \label{part_rminus1_timesS} Expand $(r-1)S$ and simplify your answer.
        
        \item \label{part_geom_sum_derive} Use your answer to (\ref{part_rminus1_timesS}) to find a formula for $S$ in terms of $r$.  Hint: Since $r \neq 1$, you may divide both sides of your equation by $r-1$.
            
        \item Use your formula from (\ref{part_geom_sum_derive}) to evaluate $S$ when $r=10$. Check that you get the same answer as in (\ref{part_sum_powers2_to4}).
            
        \item \label{part_geom_sum_ais1} Conjecture an explicit formula for $\displaystyle \sum_{k=0}^nr^k = 1 + r + r^2 + \ldots + r^n$ when $r \neq 1$.

        \item Use your formula from (\ref{part_geom_sum_ais1}) to show that $\displaystyle \sum_{k=0}^n2^k = 2^{n+1}-1$.

        \item In the Handshakes puzzle (\Cref{act_handshakes}), we saw a formula for the sum of the first $n$ integers, namely \Cref{thm_hs}, which says $$\sum_{k=1}^nk=\frac{n(n+1)}{2}.$$  Use this formula to evaluate $\displaystyle \sum_{k=1}^{10}k$.

        \item \label{part_sum_sq_formula} There is an explicit formula for the sum of the first $n$ integers squared, namely $$\sum_{k=1}^nk^2=\frac{n(n+1)(2n+1)}{6}.$$  Use this formula to evaluate $\displaystyle \sum_{k=1}^{5}k^2$.

        \item Check your answer to (\ref{part_sum_sq_formula}) by evaluating $\displaystyle \sum_{k=1}^{5}k^2$ directly.

        \item Conjecture a formula for $\displaystyle \sum_{k=1}^nk^3$ by considering the first few partial sums.  Hint:  Each partial sum is a square.  What is the pattern in the integers being squared?
    \end{apart}   
\end{act}

Make sure you complete \Cref{act_special sums} because we are about to reveal many of the answers.

\begin{exam}[Finite geometric sum conjecture]
\label{exam_finite_geom_sum}
    Consider real numbers $a$ and $r \neq 1$ and the sum
        $$S = \sum_{k=0}^6 ar^k = a + ar + ar^2 + ar^3 + ar^4 + ar^5 + ar^ 6.$$
    Note that $S$ is a function of $a$ and $r$.      
        \begin{apart}
            \item \label{part_geom_sum_exam_a3r2} Evaluate $S$ when $a=3$ and $r=2$.
                \begin{soln}
                    We get
                    $$S = \sum_{k=0}^6 3 \cdot 2^k = 3 + 3 \cdot 2 + 3 \cdot 2^2 + 3 \cdot 2^3 + 3 \cdot 2^4 + 3 \cdot 2^5 + 3 \cdot 2^ 6= 381.$$
                \end{soln}
            \item \label{part_rminus1_timesS_general} Expand $(r-1)S$ and simplify your answer.
                \begin{soln}
                    We get
                        \begin{align*}
                            (r-1)S  & = (r-1)(a + ar + ar^2 + ar^3 + ar^4 + ar^5 + ar^ 6) \\
                            & = (ar + ar^2 + ar^3 + ar^4 + ar^5 + ar^6 + ar^ 7) - (a + ar + ar^2 + ar^3 + ar^4 + ar^5 + ar^ 6) \\
                            & = \cancel{ar} + \cancel{ar^2} + \cancel{ar^3} + \cancel{ar^4} + \cancel{ar^5} + \cancel{ar^6} + ar^7 - a - \cancel{ar} - \cancel{ar^2} - \cancel{ar^3} - \cancel{ar^4} - \cancel{ar^5} - \cancel{ar^ 6}) \\
                            & = ar^7-a \\
                            & = a(r^7-1)
                        \end{align*}
                    Thus, $(r-1)S = a(r^7-1)$.
                \end{soln}            
            \item \label{part_geom_sum_derive_general} Use your answer to (\ref{part_rminus1_timesS_general}) to find a formula for $S$ in terms of $a$ and $r$. 
                \begin{soln}
                    Since $r \neq 1$, we can divide each side of the equation by $r-1$ to get $\displaystyle S = \frac{a(r^7-1)}{r-1}$.
                \end{soln}
            \item Use your formula from (\ref{part_rminus1_timesS_general}) to evaluate $S$ when $a = 3$ and $r=2$.  Check that you get the same answer as in (\ref{part_geom_sum_exam_a3r2}).
                \begin{soln}
                    We get $\displaystyle S = \frac{3(2^7-1)}{2-1}= 3(127) = 381$.
                \end{soln}
            \item \label{part_geom_sum_conj_general} Conjecture an explicit formula for $\displaystyle \sum_{k=0}^nar^k = a + ar + ar^2 + \ldots + ar^n$ when $r \neq 1$. 
                \begin{soln}
                    We conjecture that $\displaystyle \sum_{k=0}^nar^k = \frac{a(r^7-1)}{r-1} = a\frac{r^7-1}{r-1}$.
                \end{soln}
            \item Use your formula from (\ref{part_geom_sum_conj_general}) to show that $\displaystyle \sum_{k=0}^n4\cdot5^k = 5^{n+1}-1$.   
                \begin{soln}
                    We get $\displaystyle \sum_{k=0}^n4\cdot5^k = 4\frac{5^{n+1}-1}{5-1}= \cancel{4}\frac{5^{n+1}-1}{\cancel{4}} = 5^{n+1} - 1$.
                \end{soln}
        \end{apart}
\end{exam}

Let's state the general formula for a geometric sum.

\begin{thm}[Sum of finite geometric]
\label{thm_geom_sum_finite}
    For real numbers $a$ and $r \neq 1$,
        $$\sum_{k=0}^nar^k=a \frac{r^{n+1}-1}{r-1}.$$
\end{thm}

Here is an example using the finite geometric sum formula.

\begin{exam}[Using the finite geometric sum formula]
\label{exam_finite_geom_sum}
    Use the finite geometric sum formula (\Cref{thm_geom_sum_finite}) to evaluate $\displaystyle \sum_{k=0}^7(10\cdot2^k)=10 + 20 + 40 + \ldots + 1280$.
        \begin{soln}
            Our sum is geometric with $a=10$, $r=2 \neq 1$, and $n=7$ so by \Cref{thm_geom_sum_finite} we have 
            $$\sum_{k=0}^710\cdot2^k = 10 \left(\frac{2^{7+1}-1}{2-1}\right) = 10 \left(\frac{2^8-1}{1}\right) = 10(256-1) = \text{2,550}.$$
        \end{soln}
\end{exam}

We can use the distributive property (\Cref{defn_integer_algebra_order_operations}) to derive linearity properties of sums.

\begin{thm}[Linearity properties of sums]
\label{thm_linearity_sums} 
    Let $a_k$ and $b_k$ be sequences of real numbers for $k \ge s$ for some nonnegative integer $s$ and let $c$ be a real number. Then
        \begin{apart}
            \item \label{part_constant_times_sum} $\displaystyle \sum_{k=s}^n(ca_k)=c\sum_{k=s}^na_k.$  
            \item \label{part_sum_of_sum}  $\displaystyle \sum_{k=s}^n(a_k+b_k)=\sum_{k=s}^na_k+\sum_{k=s}^nb_k.$   
        \end{apart}
\end{thm}

We have seen the sum of the first $n$ integers, the handshake numbers, in many places starting with \Cref{act_handshakes}.  There are also formulas for the sum of the first $n$ squares and cubes.

\begin{thm}[Summation formulas for powers of integers]
\label{thm_sum_formula_polys}
    For a nonnegative integer $n$ we have
        \begin{apart}
            \item \label{part_sum_constant} The sum of a constant: $\displaystyle \sum_{k=1}^n1=n$.
            \item \label{part_sum_hs} The sum of the first $n$ integers: $\displaystyle \sum_{k=1}^nk = \frac{n(n+1)}{2}$.
            \item \label{part_sum_squares} The sum of the first $n$ squares: $\displaystyle \sum_{k=1}^nk^2 = \frac{n(n+1)(2n+1)}{6}$.
            \item \label{part_sum_cubes} The sum of the first $n$ cubes: $\displaystyle \sum_{k=1}^nk^3 = \left(\frac{n(n+1)}{2}\right)^2$.
        \end{apart}
\end{thm}

Here is an example using these formulas.

\begin{exam}[Using linearity properties and summation formulas]
\label{exam_sum_formulas}
   Use linearity properties and summation formulas to evaluate each sum.
       \begin{apart}
           \item Evaluate $\displaystyle \sum_{k=1}^{10}k^2=1 + 4 + 9 + \ldots + 100$
               \begin{soln}
                   By \Cref{thm_sum_formula_polys} (\ref{part_sum_squares}) with $n=10$ we have
                       $$\sum_{k=1}^{10}k^2 = \frac{10(10+1)(2(10)+1)}{6} = \frac{(10)(11)(21)}{6} = 385.$$
               \end{soln} 
           \item Evaluate $\displaystyle \sum_{k=1}^{10}k= 1 + 2 + 3 + \ldots + 10$          
               \begin{soln}
                   By \Cref{thm_sum_formula_polys} (\ref{part_sum_hs}) with $n=10$ we have
                    $$\sum_{k=1}^{10}k= \frac{10(10+1)}{2} = \frac{(10)(11)}{2} = 55.$$
               \end{soln} 
            \item Evaluate $\displaystyle \sum_{k=1}^{10}1=\underbrace{1+1+\ldots+1}_{10 \text{ times}}$
                \begin{soln}
                    By \Cref{thm_sum_formula_polys} (\ref{part_sum_constant}) with $n=10$ we have $\displaystyle \sum_{k=1}^{10}1=10.$
                    Not sure we really needed a formula here.
                \end{soln}
            \item Evaluate $\displaystyle \sum_{k=1}^{10}(3k+1)(k+2)=(4)(3)+(7)(4)+(10)(5) + \ldots + (31)(12)$.
               \begin{soln}
                   We begin by multiplying (using FOIL) to get $(3k+1)(k+2) = 3k^2+5k+2$. Then we can apply the linearity properties of sums (\Cref{thm_linearity_sums}) and our answers from the previous parts to get
                       \begin{align*}
                           \sum_{k=1}^{10}(3k+1)(k+2) 
                           & =\sum_{k=1}^{10} (3k^2+5k+2) \\
                           & = \sum_{k=1}^{10}(3k^2)+\sum_{k=1}^{10}(5k)+\sum_{k=1}^{10}2 \\
                           & = 3\left(\sum_{k=1}^{10}k^2\right)+5\left(\sum_{k=1}^{10}k\right)+2\left(\sum_{k=1}^{10}1\right) \\
                           & = 3(385) + 5(55)+2(10) = \text{1,450}.
                       \end{align*}
           \end{soln}
       \end{apart}
\end{exam}

%%%%%%%%%%%%%%%%%%%%%%%%%

\newcommand{\subfracalg}{Algebra: Fraction Arithmetic}
\subsection{\subfracalg}
\label{sub_fraction_algebra}

Some of our summation formulas involve fractions. Let's review fraction arithmetic.

\begin{defn}[Adding Fractions]
\label{defn_fraction_addition}
~
    \begin{apart}        
        \item \label{part_sum_fractions} For integers $n$, $m$, and $d$ with $d \neq 0$, the \vocab{sum of the fractions} is given by
        $$\displaystyle \frac{n}{d} + \frac{m}{d} = \frac{n+m}{d}.
        $$  For example, $$\frac{3}{7} + \frac{2}{7} = \frac{3+2}{7} = \frac{5}{7}$$  
        \item To add two fractions using (\ref{part_sum_fractions}), we need the two fractions to have the same denominator (\vocab{common denominator}).  If two fractions have different denominators, we can replace one or both fractions with equivalent fractions that have a common denominator using \Cref{thm_equal_fractions}.
    \end{apart}
\end{defn}

Let's practice working with fractions.

\begin{exam} [Adding and Subtracting Fractions]
\label{fraction arithmetic} 
    In this example, we only use computational technology to check our answers. 
    \begin{apart}
        \item Write as one fraction and simplify $\displaystyle \frac{3}{4} - \frac{1}{2}.$
            \begin{soln}
                We can use 4 as the common denominator.
                $$\displaystyle \frac{3}{4} - \frac{1}{2}
                = \frac{3}{4} - \frac{2(1)}{2(2)}
                = \frac{3}{4} - \frac{2}{4}
                = \frac{3-2}{4}
                = \frac{1}{4}
                $$
            \end{soln}
        \item Write as one fraction and simplify $\displaystyle \frac{1}{2} + \frac{2}{3}.$  
            \begin{soln}
                We can use $(2)(3) = 6$ as the common denominator.
                 $$\displaystyle \frac{1}{2} + \frac{2}{3}
                = \frac{3(1)}{3(2)} + \frac{2(2)}{2(3)}
                = \frac{3}{6} + \frac{4}{6}
                = \frac{3+4}{6}
                = \frac{7}{6}
                $$               
            \end{soln}
        \item Write as one fraction and simplify
        $\displaystyle \frac{5}{3}-2.$
            \begin{soln}
                We can rewrite 2 as the fraction $\frac{2}{1}$ and then use 3 as the common denominator.
                $$ \displaystyle \frac{5}{3}-2
                = \frac{5}{3}-\frac{2}{1}
                = \frac{5}{3}-\frac{3(2)}{3(1)}
                = \frac{5}{3}-\frac{6}{3}
                = \frac{5-6}{3}
                = \frac{-1}{3}
                =^-\frac{1}{3}
                $$
            \end{soln}
        \item Write as one fraction and simplify
        $\displaystyle \frac{1}{k-1} - \frac{1}{k}.$
            \begin{soln}
                We can use $k(k-1)$ as the common denominator.
                $$\displaystyle \frac{1}{k-1} - \frac{1}{k} 
                = \frac{k(1)}{k(k-1)} - \frac{(k-1)(1)}{(k-1)(k)} 
                = \frac{k}{k(k-1)} - \frac{k-1}{(k)(k-1)}$$
                $$ 
                = \frac{k-(k-1)}{(k)(k-1)} 
                = \frac{k-k+1}{k(k-1)} 
                = \frac{1}{k(k-1)} 
                $$

                By the way, it is sometimes useful to reverse this equation to write $$\frac{1}{k(k-1)} = \frac{1}{k-1} - \frac{1}{k}.$$  
            \end{soln}

        \item Write as one fraction and simplify
        $\displaystyle \frac{1}{2^{k-1}} + \frac{1}{2^k}.$
            \begin{soln}
                Recall that $2\cdot 2^{k-1} =2^k$ so we can 
                use $2^k$ as the common denominator.
                $$\displaystyle \frac{1}{2^{k-1}} + \frac{1}{2^k}
                = \frac{2(1)}{2(2^{k-1})} + \frac{1}{2^k}
                = \frac{2}{2^k} + \frac{1}{2^k}
                = \frac{2+1}{2^k}
                = \frac{3}{2^k}
                $$
            \end{soln}
        \item Write as one fraction and simplify
        $\displaystyle \frac{1}{(n-1)!} + \frac{1}{n!}.$
            \begin{soln}
                Recall that $n!=n(n-1)!$ so we can use $n!$ as the common denominator to get
                    $$\displaystyle \frac{1}{(n-1)!} + \frac{1}{n!}
                    = \frac{n(1)}{n(n-1)!)} + \frac{1}{n!}
                    = \frac{n}{n!} + \frac{1}{n!}
                    = \frac{n+1}{n!}$$
             \end{soln}
    \end{apart}
\end{exam}

Adding or subtracting fractions is a bit complicated because we need to find a common denominator.  Luckily, multiplying fractions is simpler.

\begin{defn}[Multiplying Fractions]
\label{defn_multiplying_fractions}
~
    \begin{apart}
        \item For integers $n, m, c, d$ with $c,d \neq 0$, the \vocab{product} of the fractions is given by
        $$\frac{m}{c} \cdot \frac{n}{d} = \frac{mn}{cd}.$$
        For example, 
        $$\frac{2}{3} \cdot \frac{7}{5} 
        = \frac{(2)(7)}{(3)(5)} = \frac{14}{15}.$$
    \end{apart}
\end{defn}

Here is an example.

\begin{exam}[Multiplying Fractions]
\label{exam_multiplying_fractions}
In this example, we only use computational technology to check our answers.
    \begin{apart}
        \item Write as one fraction and simplify: $\displaystyle \frac{1}{2} \cdot \frac{5}{6}.$
            \begin{soln}
                By \Cref{defn_multiplying_fractions},
                $$\frac{1}{2} \cdot \frac{5}{6}
                = \frac{(1)(5)}{(2)(6)}
                = \frac{5}{12}.
                $$
            \end{soln}
        \item Write as one fraction and simplify $\displaystyle \frac{1}{2} \cdot \frac{2}{3}$.
            \begin{soln}
                We have $$\frac{1}{2} \cdot \frac{2}{3} = \frac{(1)(2)}{(2)(3)}= \frac{\cancel{2}(1)}{\cancel{2}(3)}=\frac{1}{3}.$$
            \end{soln}
        \item Write as one fraction and simplify $\displaystyle 5 \cdot \frac{3}{10}$.
            \begin{soln}
                To use \Cref{defn_multiplying_fractions}, we need to write 5 as the fraction $\frac{5}{1}$.  Next,
                $$5 \cdot \frac{3}{10}
                = \frac{5}{1} \cdot \frac{3}{10}
                = \frac{(5)(3)}{(1)(10)}
                = \frac{\cancel{5}(3)}{\cancel{5}(2)}
                = \frac{3}{2}.
                $$
            \end{soln}
        \item Simplify $\displaystyle 4 \left(\frac{4^m-1}{3}\right)+1$
            \begin{soln}
                First, we can write 4 as $\frac{4}{1}$ and multiply the fractions to get
                $$4 \left(\frac{4^m-1}{3}\right)+1
                = \frac{4}{1} \left(\frac{4^m-1}{3}\right)+1
                = \frac{4\left(4^m-1\right)}{3}+1
                .$$
                Next, we can use 3 as the common denominator and write 1 as $\frac{3}{3}$ to get
                $$\frac{4\left(4^m-1\right)}{3}+1
                = \frac{4\left(4^m-1\right)}{3}+\frac{3}{3}
                = \frac{4\left(4^m-1\right)+3}{3}
                .$$
                Next, using the distributive property and simplifying we get
                $$\frac{4\left(4^m-1\right)+3}{3}
                = \frac{4\left(4^m\right)-4(1)+3}{3}
                =\frac{4\left(4^m\right)-1}{3}       
                .$$  
                Last, we can simplify the power of 4 to get
                $$\frac{4\left(4^m\right)-1}{3} 
                = \frac{4^{m+1}-1}{3}$$
            \end{soln}
    \end{apart}
\end{exam}

It is your turn to practice working with fractions.

\begin{act} [Fractions] 
    \label{act_fractions}  
    On this problem, you may only use technology to check your answers.
    \begin{actpart} 
        \item Write as one fraction and simplify $\displaystyle \frac{3}{10} - \frac{4}{5}$.
        \item Write as one fraction and simplify $\displaystyle \frac{3}{10} \cdot \frac{5}{6}$.
        \item Write as one fraction $\displaystyle \frac{1}{n}+ \frac{1}{n^2}$.
        \item Write as one fraction and simplify $\displaystyle \frac{1}{2^{k-1}} - \frac{1}{2^k}$
        \item Write as one fraction and simplify $\displaystyle 3 \left(\frac{3^m-1}{2}\right)-1$
    \end{actpart}
\end{act}

%%%%%%%%%%%%%%%%%%%%%%%%

\subsection{Exercises}

%%%%%%%%%%%%%%%%%%%%%%%%

\subsubsection*{Exercises for \Cref{sub_pff_sum_proofs_induction}
\subpffsumproofsind}

\begin{exer}[\exerP] %  sum of odds]
\label{exer_math_ind_sum_odds}
    Use mathematical induction to prove that $\displaystyle \sum_{k=1}^n(2k -1)=n^2$ for $n\ge 1$.
\end{exer}

\begin{exer}[\exerP] %  sum of evens]
\label{exer_math_ind_sum_evens}
    Use mathematical induction to prove that $\displaystyle \sum_{j=1}^n(2j) = n(n+1)$ for $n \ge 1$.
\end{exer}

\begin{exer}[\exerP] %  sum of factorials]
\label{exer_math_ind_sum_factorials}
    Use mathematical induction to prove that $\displaystyle \sum_{k=1}^nk(k!) = (n+1)!-1$
    for $n\ge 1$. We saw this sum in \Cref{act_fact_rep}.
\end{exer}

\begin{exer}[\exerP] %  geom sum with powers of 5]
\label{exer_math_ind_geom_sum5}
    Use mathematical induction to prove that $\displaystyle \sum_{k=1}^n4 \cdot 5^k = 5^{n+1}-1$ for $n\ge 1$.
\end{exer}

\begin{exer}[\exerP] %  prove product 1 plus reciprocal]
\label{exer_math_ind_prod_1plus_recip}
    Use mathematical induction to prove that $\displaystyle \prod_{k=1}^n \left(1 + \frac{1}{k}\right) = n+1$ for $n\ge 1$.
\end{exer}

\begin{exer}[\exerP] %  prove product evens]
\label{exer_math_ind_prod_evens}
    Use mathematical induction to prove that $\displaystyle \prod_{j=1}^n(2j) = 2^nn!$ for $n \ge 1$.
\end{exer}

\begin{exer} [\exerP] %  sum first $n$ Fibonaccis proof]  
\label{exer_sum_firstn_Fib_proof}
    Use mathematical induction to prove that $\displaystyle\sum_{k=1}^nF_k = F_{n+2}-1$ for $n \ge 1$ where $F_n$ are the Fibonacci numbers.   % From  \Cref{sec_sum_conj} : \Cref{exer_sum_firstn_Fib_conj}
\end{exer}

\begin{exer}[\exerU] %  prove sum $4k+1$]
\label{exer_math_ind_sum_4kplus1}
    Use mathematical induction to prove that $\displaystyle \sum_{k=1}^n(4k-1) = n(2n+1)$ for $n\ge 0$.   
\end{exer}

\begin{exer}[\exerU] %  prove sum reciprocals power of 2]
\label{exer_math_ind_sum_recip_powers2}
   Use mathematical induction to prove that $\displaystyle \sum_{k=0}^n\frac{1}{2^k} = 2 - \frac{1}{2^n}$ for $n \ge 1$.
\end{exer}

\begin{exer}[\exerU] %  sum first $n$ Fibonaccis squared proof]
\label{exer_sum_firstn_Fib_sq_proof}
    Use mathematical induction to prove that $\displaystyle\sum_{k=1}^nF_k^2= F_nF_{n+1}$ for $n \ge 1$ where $F_n$ are the Fibonacci numbers. % From  \Cref{sec_sum_conj} : \Cref{exer_sum_firstn_Fib_sq_conj} 
\end{exer}

\begin{exer}[\exerU] % Alternating sum Fibonacci
\label{exer_sum_alt_fib_proof}
    Use mathematical induction to prove that    
    $\displaystyle\sum_{k=1}^n(-1)^{k+1}F_k = 1+(-1)^{n+1}F_{n-1}$ for $n \ge 1$ where $F_n$ are the Fibonacci numbers. % From  \Cref{sec_sum_conj} : \Cref{act_fib_sums_conj}  
\end{exer}

\begin{exer}[\exerU] %  prove sum $k\cdot 2^k)$]
\label{exer_math_ind_sum_multpowers2}
    Use mathematical induction to prove that $\displaystyle \sum_{k=1}^nk \cdot 2^k = (n-1)2^{n+1}+2$ for $n \ge  1$.   
\end{exer}

\begin{exer}[\exerU] %  prove sum alternating squares]
\label{exer_math_ind_sum_alt_sq}
    Use mathematical induction to prove that 
    
    $\displaystyle  \sum_{k=0}^n(-1)^kk^2 = (-1)^n ~ \frac{n(n+1)}{2}$
    for $n\ge 0$. 
\end{exer}

\begin{exer}[\exerU] %  prove \Cref{thm_sum_formula_polys} (\ref{part_sum_squares})]
\label{exer_math_ind_sum_squares}
    Use mathematical induction to prove that $\displaystyle \sum_{k=1}^nk^2=\frac{n(n+1)(2n+1)}{6}$ for $n\ge 1$.  This exercise proves \Cref{thm_sum_formula_polys} (\ref{part_sum_squares}).
\end{exer}

\begin{exer}[\exerU] %  prove \Cref{thm_sum_formula_polys} (\ref{part_sum_cubes})]
\label{exer_math_ind_sum_squares}
    Use mathematical induction to prove that $\displaystyle \sum_{k=1}^nk^3=\left(\frac{n(n+1)}{2}\right)^2$ for $n\ge 1$. This exercise proves \Cref{thm_sum_formula_polys} (\ref{part_sum_cubes}).
\end{exer}

\begin{exer}[\exerU] %  prove sum $(3k+1)(k+2)$]
\label{exer_math_ind_sum_poly}
    Use mathematical induction to prove that 
    
    $\displaystyle \sum_{k=1}^n(3k+1)(k+2) = n(n+2)(n+3)$ for $n\ge 1$. 
\end{exer}

\begin{exer}[\exerU] %  prove sum fraction]
\label{exer_math_ind_sum_frac}
    Use mathematical induction to prove that 
    
    $\displaystyle \sum_{k=1}^n\frac{1}{(2k-1)(2k+1)}=\frac{n}{2n+1}$ for $n\ge 1$.    
\end{exer}

\begin{exer}[\exerU] %  prove sum reciprocal handshakes]
\label{exer_math_ind_sum_hs_recip}
    Use mathematical induction to prove that     
    $\displaystyle \sum_{k=1}^n\frac{2}{k(k+1)}=\frac{2n}{n+1}$ for $n\ge 1$.    
\end{exer}

\begin{exer}[\exerR] % sum and product proofs by mathematical induction
\label{exer_dyk_sum_proofs}
    Do you know \ldots
        \begin{exerpart}
            \item What the recursive definitions of partial sums and partial products are?
            \item How to use proof by mathematical induction to verify the explicit formula of a sum or product?
        \end{exerpart}
\end{exer}

\begin{exer}[\exerE] %  prove \Cref{thm_geom_sum_finite}]
\label{exer_math_ind_geom_sum_formula}
    Use mathematical induction to prove  that
    
    $\displaystyle \sum_{k=0}^nar^k=a\left(\frac{r^{n+1}-1}{r-1}\right)$ for $n \ge 0$. This exercise proves \Cref{thm_geom_sum_finite}.
\end{exer}

%%%%%%%%%%%%%%%%%%%%%%%%

\subsubsection*{Exercises for \Cref{sub_pff_comb_sums} \subpffcombsums}

\begin{exer}[\exerP] %  counting cats and dogs]
\label{exer_count_cats_dogs}
    Suppose there are $n$ cats and $n$ dogs and we want to choose $n$ animals.
        \begin{exerpart}
            \item If we picked $k$ cats, then how many dogs do we need to choose?
            \item How many ways are there to pick $k$ cats and the rest dogs?
            \item \label{part_picked_fav_cat} If we picked one favorite animal, how many ways are there to pick the rest?  Hint: how many animals are left and how many more do we need?
        \end{exerpart}
\end{exer}

\begin{exer}[\exerU] %  combinatorial proof of the sum of the binomial coefficients squared]
\label{exer_comb_pf_sum_binom_sq}
    Give a combinatorial proof that 
        $$\displaystyle \sum_{k=0}^n\binom{n}{k}^2 = \binom{2n}{n}$$ 
    for $n \ge 1.$ Hint: Suppose there are $n$ cats and $n$ dogs and we want to choose $n$ animals.  Condition on $k=$ the number of cats chosen.  
\end{exer}

\begin{exer}[\exerU] %  combinatorial proof of the weighted  sum of the binomial coefficients squared]
\label{exer_comb_pf_sum_weighted_binom_sq}
  Give a combinatorial proof that 
      $$ \sum_{k=1}^nk\binom{n}{k}^2 = n\binom{2n-1}{n-1}$$ 
    for $n \ge 1$. Hint:  Suppose there are $n$ cats and $n$ dogs and we want to choose $n$ animals including a favorite cat.  Condition on $k=$ the number of cats chosen.  
\end{exer}

\begin{exer}[\exerR] % combinatorial proof of sum
\label{exer_dyk_comb_pf_sums}
    Do you know \ldots
        \begin{exerpart}
            \item How to write a combinatorial proof of a sum?
            \item What the phrase ``condition on $k$'' indicates?
        \end{exerpart}
\end{exer}

\begin{exer}[\exerE] %  combinatorial proof of handshakes formula]
\label{exer_comb_pf_sum_hs}
    Give a combinatorial proof that 
        $$\sum_{k=1}^nk = \binom{n+1}{2}$$ 
    for $n \ge 1$. Hint:  How many 2-elements subsets does $S = \{0, 1, 2, \ldots, n\}$ have?  Condition on $k =$ the larger number in the subset.  Note that the other number in the subset must come from $\{0, 1, 2, \ldots, k-1\}$.  
\end{exer}

\begin{exer}[\exerE] %  combinatorial proof of Hockey Stick]
\label{exer_comb_pf_sum_hockey_stick}
    Give a combinatorial proof (using \Cref{pff_comb_proof_sum}) that  
        $$\sum_{k=m}^n \binom{k}{m} = \binom{n+1}{m+1}$$
    for $m, n \ge 1$. Hint:  How many $(m+1)$-element subsets does $S = \{0, 1, 2, \ldots, n\}$ have?  Condition on $k =$ the largest number in the subset.  Note that the other $m$ numbers in the subset must come from $\{0, 1, 2, \ldots, k-1\}$.  This proves the hockey stick conjecture from \Cref{exam_hockey_stick}.
\end{exer}

\begin{exer}[\exerE] % Comb proof factorial id
\label{exer_comb_pf_factorials}
    Consider the set $S = \{0, 1, 2, \ldots, n\}$.
    \begin{exerpart}
        \item \label{part_counting_nonid_perms} How many nonidentity permutations of $S$ are there?  Hint: $|S|=n+1$.
        \item \label{part_counting_perms_that_fix_afterk} Consider an integer $k$ with $0 \le k \le n$.  How many permutations of $S$ are there where $k$ is the largest integer that is not fixed?  
        
        Hint: Such a permutation has the format
            $$\sigma = \left(\begin{array} {ccccccccccc}  
            0 & 1 & 2 & \cdots & k-1 & k & k+1 & k+2 & \cdots & n-1 & n\\
            \underline{~\quad} & \underline{~\quad} & \underline{~\quad} & \cdots &  \underline{~\quad} & \underline{~\quad}
            & k+1 & k+2 &  \cdots & n-1 & n  \end{array}\right)$$
        Since $\sigma(k) \neq k$, the choices for $\sigma(k)$ are $0, 1, 2, \ldots, k-1$.  That fills in the blank under $k$ and leaves $k$ integers to be arranged into the remaining $k$ blanks.
        \item Give a combinatorial proof that $\displaystyle \sum_{k=1}^nk\cdot k! = (n+1)!-1$ for $n \ge 1$. 
        
        Hint:  Use part (a) for one count.  For the other count, condition on $k=$ the largest integer not fixed by the nonidentity permutation and use part (b).
    \end{exerpart}
\end{exer}

\begin{exer}[\exerE] %Fibonacci sums in the arithmetic triangle
\label{exer_fib_in_pascal}
    In \Cref{act_squares_dominoes}, we saw that the number of ways to tile a $1 \times n$ board using $1 \times 1$ squares and $1 \times 2$ dominoes was $F_{n+1}$, where $F_n$ are the Fibonacci numbers.  In this exercise, we use this characterization of the Fibonacci numbers to give a combinatorial proof that $$\sum_{k = 0}^n \binom{n-k}{k} = F_{n+1}.$$
        \begin{exerpart}
            \item \label{part_sq_as_fn_dom} We want to tile the $1 \times n$ board using $s$ squares and $k$ dominoes. Note that there are $k+s$ tiles. In how many ways can we tile the board?  Hint: Think of $k+s$ places to put tiles.  Choose which $k$ places will be dominoes. % Ans \binom{k+s}{k} 
            \item Let's rewrite the answer to (\ref{part_sq_as_fn_dom}) using $k$ and $n$ (instead of $s$).  First, explain why $2k+s=n$.  Next, solve for $s$. Last, substitute into your answer to (\ref{part_sq_as_fn_dom}) to get a count of the number of ways to tile the board using $k$ dominoes in terms of $k$ and $n$.
            \item Count the total number of ways to tile the $1 \times n$ using squares and dominoes by using cases conditioned on $k=$ the number of dominoes.  Since cases add (\Cref{thm_cases_add}), your answer should be a summation.
\end{exerpart}
\end{exer}

%%%%%%%%%%%%%%%%%%%%%%%%

\subsubsection*{Exercises for \Cref{sub_special_sums} \subspecialsum}

\begin{exer}[\exerP] %  evaluate the sum of a finite geometric using formulas]
\label{exer_eval_sum_formulas_geom}
    Use \Cref{thm_geom_sum_finite} to evaluate each sum.
        \begin{exerpart}
            \item Evaluate $\displaystyle \sum_{k=0}^{8}2 \cdot 3^k$.
            \item Evaluate $\displaystyle \sum_{k=0}^{8}10 \cdot (-1)^k$.
        \end{exerpart}
\end{exer}

\begin{exer}[\exerP] %  evaluate the sum of polynomials using formulas]
\label{exer_eval_sum_formulas_polys}
    Use \Cref{thm_sum_formula_polys} to evaluate each sum.
        \begin{exerpart}
            \item Evaluate $\displaystyle \sum_{k=1}^{8}k$.
            \item Evaluate $\displaystyle \sum_{k=1}^{8}k^2$.
            \item Evaluate $\displaystyle \sum_{k=1}^{8}k^3$.
        \end{exerpart}
\end{exer}

\begin{exer}[\exerU] % practice using polys and lienarity
\label{exer_using_sum_formulas_polys_lin}
     Use \Cref{thm_linearity_sums} and \Cref{thm_sum_formula_polys} to evaluate $\displaystyle \sum_{k=1}^n(2k+1)(k+3)$ and write your answer in factored form.  
\end{exer}

\begin{exer}[\exerU] %  sum and product of powers of $x$]
\label{exer_sum_prod_powersofx}
    Use summation formulas to simplify each of the following expressions. (Yes, each uses summation formulas.)  
        \begin{exerpart}
            \item $1 + x + x^2 + \ldots + x^n$
            \item $xx^2x^3...x^n$
        \end{exerpart}
\end{exer}

\begin{exer}[\exerR] % special sums
\label{exer_dyk_special_sums}
    Do you know \ldots
        \begin{exerpart}
            \item What the formula is for a finite geometric sum and how to use it?
            \item What the linearity properties of sums say?
            \item How to use the formulas for the sum of the first $n$ integers, their squares, and their cubes?
        \end{exerpart}
\end{exer}

\begin{exer}[\exerE] %  perfect integers]
\label{exer_perfect}
    The integer $b$ is a \vocab{proper divisor} of $n$ if $b \mid n$ but $b < n$. An integer is \vocab{perfect} if it equals the sum of its proper divisors. 
        \begin{exerpart}
            \item Verify that the integer 28 is perfect.
            \item What is the smallest perfect integer?  Justify your answer. That means justify that it's perfect and also that it's smallest.
            \item Can a prime number be perfect?  Explain.
            \item \label{part_geom_sum_perfect} Show how to use the finite geometric sum formula (\Cref{thm_geom_sum_finite}) to simplify $\displaystyle \sum_{k=0}^{p-1}2^k$ where $p$ is a positive integer.
            \item \label{part_divisors_power2_times_odd_prime} Suppose $n = 2^{p-1}q$ where $q$ is an odd prime.  List the divisors of $n$, using $\ldots$ as needed.  Note that $2^{p-1}q$ is the prime factorization of $n$.  
            \item  \label{part_sum_divisors_perfect} Calculate the sum of the divisors of $n$ from (\ref{part_divisors_power2_times_odd_prime}) and use your answer from (\ref{part_geom_sum_perfect}) to simplify your answer.
            \item Note that $n$ is perfect if the sum of \emph{all} of the divisors of $n$ is $n+n=2n$. Suppose $p$ is prime and $q=2^p-1$ also happens to be prime.  Observe that $q$ is odd. Consider the integer $n = 2^{p-1}q$ as before. Use (\ref{part_sum_divisors_perfect}) to prove that $n$ is perfect.
        \end{exerpart}
\end{exer}

\begin{exer}[\exerE] %  infinite geometric sum]
\label{exer_infinite_geom_sum}
    For readers who have studied calculus, an infinite series is defined as a limit.  For example,
       
 $$\sum_{k=1}^\infty (ar^k) = \lim_{n \to \infty} \sum_{k=1}^n (ar^k)$$
    \begin{exerpart}
        \item Use the finite geometric sum formula (\Cref{thm_geom_sum_finite}) to simplify $\displaystyle \sum_{k=1}^n (ar^k)$. 
        \item If $-1<r<1$, explain why $\displaystyle \lim_{n \to \infty} r^{n+1} = 0$.
        \item Evaluate $\displaystyle \sum_{k=1}^\infty (ar^k)$  when $-1 < r < 1$.  Simplify your answer, including multiplying the top and bottom of the your limit by -1 (using \Cref{thm_equal_fractions}).
    \end{exerpart}
\end{exer}

%%%%%%%%%%%%%%%%%%%%%%

\subsubsection*{Exercises for \Cref{sub_fraction_algebra}
\subfracalg}

\begin{exer}[\exerP] %  adding and subtracting fractions]
\label{exer_add_sub_fractions}
    On this problem, do not use computational technology except to check.
        \begin{exerpart}
            \item Write as one fraction and simplify: $\displaystyle \frac{3}{4} + \frac{1}{2}$.
            \item Write as one fraction and simplify: $\displaystyle \frac{3}{4} \cdot \frac{2}{7}$.
        \end{exerpart}
\end{exer}

\begin{exer}[\exerU] %  adding and subtracting fractions with $n$]
\label{exer_add_sub_fractions_withn}
~
    \begin{exerpart}
        \item Write as one fraction and simplify: $\displaystyle \frac{n-1}{2n-1}+\frac{1}{(2n-1)(2n+1)}$.
        \item Write as one fraction and simplify: $\displaystyle \frac{2n-2}{n}+\frac{2}{n(n+1)}$.
    \end{exerpart}
\end{exer}

\begin{exer}[\exerR] % add mult fractions
\label{exer_dyk_fraction_arith}
    Do you know \ldots
        \begin{exerpart}
            \item How to add and subtract fractions?
            \item Why we need a common divisor to add and subtract fractions?
            \item How to multiply fractions?
        \end{exerpart}
\end{exer}

%%%%%%%%%%%%%%%%%%%%%%%%%%%
\subsection{\answers}
%%%%%%%%%%%%%%%%%%%%%%%%%%%

\subsubsection{\answers for \Cref{sub_pff_sum_proofs_induction}
\subpffsumproofsind}

\subsubsection*{\Cref{exer_math_ind_sum_odds} \answers}
Hints:  Make sure you use parentheses in $f(n-1)=(n-1)^2$.

\subsubsection*{\Cref{exer_math_ind_sum_factorials} \answers}
Hint: To simplify $n!-1+nn!$, first put the factorials together and factor to get $$n!+nn!-1 = (1+n)n!-1.$$  Now, simplify $(1+n)n!$ using $(1+n)=(n+1)$.

\subsubsection*{\Cref{exer_math_ind_geom_sum5} \answers}
Hint:  To simplify $5^n-1+4\cdot5^n$, first put the powers of five together and factor to get $$5^n+4\cdot5^n-1 = 5^n(1+4)-1 = 5^n\cdot5 - 1.$$

\subsubsection*{\Cref{exer_math_ind_prod_1plus_recip} \answers}
Hint: Note that $\displaystyle 1 + \frac{1}{n}  = \frac{n}{n} + \frac{1}{n}  =\frac{n+1}{n}$.

\subsubsection*{\Cref{exer_math_ind_prod_evens} \answers}
Hint: Note that $f(n-1) = 2^{n-1}(n-1)!$.  Simplify $\left(2^{n-1}(n-1)!\right)(2n)$ by collecting the powers of two and separately collecting the factorial.

\subsubsection*{\Cref{exer_sum_firstn_Fib_proof} \answers}
Hint: Note that $f(n-1) = F_{n+1}-1$ and that the Fibonacci sequence includes $$\ldots, F_n, F_{n+1}, F_{n+2}, \ldots$$

\subsubsection*{\Cref{exer_math_ind_sum_4kplus1} \answers}
Hint: Use two sets of parentheses to evaluate $f(n-1)$.

\subsubsection*{\Cref{exer_sum_firstn_Fib_sq_proof} \answers}
Hints: Note that $$F_{n-1}F_n + F_n^2 = F_n(F_{n-1} + F_n)$$

\subsubsection*{\Cref{exer_math_ind_sum_multpowers2} \answers}
Hint: Make sure to use parentheses when evaluating $f(n-1)$.  Also, note that $$n2^n + n2^n = 2n2^n=n(2\cdot2^n)$$
and $$n2^{n+1}-2^{n+1}=(n-1)2^{n+1}.$$

\subsubsection{\answers for \Cref{sub_pff_comb_sums} \subpffcombsums}

\subsubsection*{\Cref{exer_count_cats_dogs} \answers}
\begin{apart}
    \item $n-k$
    \item Hint: Step 1 is to pick $k$ out of $n$ cats.  Step 2 is to pick $n-k$ out of $n$ dogs.  Steps multiply.
    \item Hint: There are $n+n=2n$ animals to start, so after we pick a favorite animal there are $2n-1$ animals remaining.
\end{apart}

\subsubsection*{\Cref{exer_comb_pf_sum_binom_sq} \answers}
Hints: Make sure to choose $n$ animals in total.  If you choose $k$ cats, then you will need to choose $n-k$ dogs.  Look at \Cref{thm_sym_arith_triangle}.

\subsubsection*{\Cref{exer_comb_pf_sum_weighted_binom_sq} \answers}
Hints: Make sure to choose $n$ animals in total.  If you choose $k$ cats, then you will need to choose $n-k$ dogs. Also, look at \Cref{exer_count_cats_dogs} (\ref{part_picked_fav_cat}).
\subsubsection{\answers for \Cref{sub_special_sums} \subspecialsum}

\subsubsection*{\Cref{exer_eval_sum_formulas_geom} \answers}
\begin{apart}
    \item Hint: The answer is just less than 20,000.
    \item 10
\end{apart}

\subsubsection*{\Cref{exer_eval_sum_formulas_polys} \answers}
\begin{apart}
    \item Hint: The answer is less than 50.
    \item Hint: The answer is just over 200.
    \item Hint: The answer is just under 1,300.
\end{apart}

\subsubsection*{\Cref{exer_sum_prod_powersofx} \answers}
\begin{apart}
    \item Hint: Use \Cref{thm_geom_sum_finite}.
    \item Hint: Write the product as a single power of $x$ where the exponent is a sum written in sigma notation.
\end{apart}















